\chapter{Flux de détail et wholesaling}\label{ch:m1:retail-flow-and-wholesaling}

On ouvre une application de courtage, on appuie sur «~acheter~», et dix actions apparaissent
sur son compte un quart de seconde plus tard, à un prix meilleur d'une fraction de centime
que celui affiché à l'écran. On n'a payé aucune commission. L'ordre n'a jamais atteint
une bourse~: le courtier l'a transmis à une société de trading, qui a vendu les actions
depuis son stock propre et a versé au courtier une rémunération pour ce service. Les dealers
qui négocient en principal hors des bourses représentent environ quatre actions sur dix
échangées aux États-Unis
(\cref{dat:m1:us-equity-market-structure:where}), et les ordres de détail traités de cette façon
en constituent une part importante. Ce chapitre explique pourquoi une société de trading
paie pour négocier contre le client, comment on mesure la qualité de ce qu'il a reçu
et ce que ce dispositif change pour tous les autres.

\section{Le parcours d'un ordre de détail}

\begin{definition}[Wholesaler et internalisation]\label{def:m1:retail-flow-and-wholesaling:wholesaler}
Un \emph{wholesaler}\index{wholesaler} est une société de \omterm{def:m1:what-a-trading-firm-does:market-maker}{teneur de marché} qui reçoit
les ordres exécutables des clients des courtiers de détail et en exécute la plupart
en principal, pour son propre compte, en dehors de toute bourse.
L'\emph{internalisation}\index{internalisation} est l'exécution d'un ordre client
contre le stock propre de la société d'exécution plutôt que contre les ordres
d'autres investisseurs sur une plateforme.
\end{definition}

\begin{omfigure}
\begin{tikzpicture}[font=\small,
  box/.style={draw=omDef,rounded corners=2pt,minimum width=2.6cm,minimum height=0.95cm,align=center,fill=omDef!4},
  ext/.style={draw=omExo,rounded corners=2pt,minimum width=2.6cm,minimum height=0.95cm,align=center,fill=omExo!6},
  flow/.style={-{Stealth[length=2mm]},thick,omDef},
  pay/.style={-{Stealth[length=2mm]},dashed,thick,omThm}]
\node[ext] (c) at (0,0) {Client\\de détail};
\node[box] (b) at (4.3,0) {Courtier de détail\\(agent)};
\node[box] (w) at (8.6,0) {Wholesaler\\(en principal)};
\node[ext] (x) at (12.9,1.3) {Bourses};
\node[ext] (d) at (12.9,-1.3) {Dark pools};
\draw[flow] (c) -- node[above=10pt,font=\footnotesize]{ordre} (b);
\draw[flow] (b) -- node[above=10pt,font=\footnotesize]{acheminé} (w);
\draw[flow] (w) -- (x);
\draw[flow] (w) -- (d);
\draw[pay] (w.south west) to[bend left=28] node[below,font=\footnotesize]{paiement pour flux d'ordres} (b.south east);
\draw[pay] (b.south west) to[bend left=28] node[below,font=\footnotesize]{commission nulle} (c.south east);
\end{tikzpicture}

\omcaption{Un ordre au marché de détail. Le \omterm{def:m1:retail-flow-and-wholesaling:wholesaler}{wholesaler} le remplit à partir de son stock au
meilleur prix national ou légèrement mieux, verse une rémunération au courtier, et gère
la position qui en résulte sur les bourses et dans les \omterm{def:m1:us-equity-market-structure:dark}{dark pools}. Seule cette
couverture résiduelle atteint un carnet d'ordres public.}\label{fig:m1:retail-flow-and-wholesaling:path}
\end{omfigure}

Le courtier reste un agent~: il doit à son client la meilleure exécution, et le choix du
\omterm{def:m1:retail-flow-and-wholesaling:wholesaler}{wholesaler} relève de sa décision d'acheminement. Le \omterm{def:m1:retail-flow-and-wholesaling:wholesaler}{wholesaler} est un principal~:
il ne doit au client rien d'autre que le prix qu'il a promis au courtier
de livrer, et se dispute cette promesse avec d'autres \omterm{def:m1:retail-flow-and-wholesaling:wholesaler}{wholesalers}.

\begin{definition}[Paiement pour flux d'ordres et amélioration de prix]\label{def:m1:retail-flow-and-wholesaling:pfof}
Le \emph{paiement pour flux d'ordres}\index{paiement pour flux d'ordres} (PFOF) est une
rémunération versée par une société d'exécution à un courtier, en échange de l'acheminement
des ordres de ses clients vers cette société. L'\emph{amélioration de prix}\index{amélioration de prix}
est l'écart de prix dont l'exécution est meilleure que le meilleur
cours national affiché du côté client à l'instant considéré~: pour un achat, le meilleur cours vendeur
national, diminué du prix payé.
\end{definition}

\begin{dated}{2026-09}{L'ampleur du dispositif et les rapports qui le concernent}\label{dat:m1:retail-flow-and-wholesaling:scale}
Selon une compilation de rapports publics d'acheminement, réalisée par une publication spécialisée, les
courtiers de détail américains ont reçu environ 950~millions de dollars de \omterm{def:m1:retail-flow-and-wholesaling:pfof}{paiements pour flux d'ordres} au
seul deuxième trimestre 2025, pour les actions et les options réunies. Les modifications
de la règle d'exécution de la SEC (Rule 605), adoptées le 6 mars 2024, étendent
la publication mensuelle aux \omterm{def:m1:the-sell-side:sell-side}{broker-dealers} de grande taille et ajoutent un rapport
de synthèse simplifié~; leur date de mise en conformité a été reportée du 14 décembre 2025
au 1\up{er}~août 2026, les statistiques d'\omterm{def:m1:retail-flow-and-wholesaling:pfof}{amélioration de prix} par rapport au meilleur cours
affiché devant suivre en novembre 2026.
\end{dated}

\section{Mesurer ce que le client a reçu}

\begin{definition}[Demi-fourchette effective et demi-fourchette réalisée]\label{def:m1:retail-flow-and-wholesaling:effective}
Pour un ordre de sens $\epsilon = \pm1$ exécuté au prix $p$ lorsque le milieu
était $m_t$, la \emph{demi-fourchette effective}\index{demi-fourchette effective} est
$\epsilon(p - m_t)$~: ce que l'ordre a payé par rapport au milieu. La
\emph{demi-fourchette réalisée}\index{demi-fourchette réalisée} à l'horizon $\tau$ est
$\epsilon(p - m_{t+\tau})$~: ce que le fournisseur de liquidité a conservé une fois
que le prix a bougé. Ces deux grandeurs sont en général publiées doublées, comme des
fourchettes, et exprimées en centimes ou en points de base.
\end{definition}

\begin{proposition}[L'identité de la fourchette]\label{prop:m1:retail-flow-and-wholesaling:identity}
Pour toute transaction et tout horizon,
\[
\underbrace{\epsilon(p - m_t)}_{\text{effective}}
\;=\;
\underbrace{\epsilon(p - m_{t+\tau})}_{\text{réalisée}}
\;+\;
\underbrace{\epsilon(m_{t+\tau} - m_t)}_{\text{impact sur le prix}} .
\]
Moyennée sur un flux d'ordres, le terme d'impact sur le prix correspond à la \omterm{def:m1:what-a-trading-firm-does:adverse-selection}{sélection
adverse} que ce flux fait subir (\cref{def:m1:what-a-trading-firm-does:adverse-selection}),
et la \omterm{def:m1:retail-flow-and-wholesaling:effective}{demi-fourchette réalisée} représente le revenu brut par action du fournisseur de liquidité.
\end{proposition}
\begin{proof}
Ajouter et retrancher $\epsilon\,m_{t+\tau}$.
\end{proof}

L'identité fait passer l'idée centrale du \cref{ch:m1:what-a-trading-firm-does}
à une forme mesurable à partir d'un fichier de transactions et de cotations. Elle montre
aussi ce qu'une statistique de qualité d'exécution peut et ne peut pas dire~: la \omterm{def:m1:retail-flow-and-wholesaling:effective}{demi-fourchette
effective} est ce que le \emph{client} a payé~; la répartition entre les deux termes de droite
décrit le \emph{flux}, et non la qualité de l'exécution.

\begin{definition}[Rapports Rule 605 et Rule 606]\label{def:m1:retail-flow-and-wholesaling:reports}
Un \emph{rapport Rule 605}\index{rapport Rule 605} est la publication mensuelle
publique par laquelle une plateforme de négociation américaine indique, par titre, par type d'ordre et par taille,
sa vitesse d'exécution, ses fourchettes effectives et réalisées ainsi que son \omterm{def:m1:retail-flow-and-wholesaling:pfof}{amélioration de
prix}. Un \emph{rapport Rule 606}\index{rapport Rule 606} est la publication trimestrielle
publique par laquelle un courtier indique vers quelles plateformes il a acheminé
les ordres de ses clients et quels paiements il a perçus de chacune.
\end{definition}

\begin{example}[Lecture d'une exécution]\label{ex:m1:retail-flow-and-wholesaling:onefill}
La meilleure cotation nationale est $20.00 \times 20.02$~; un achat au marché
de 100 actions est exécuté à 20.018. \omterm{def:m1:retail-flow-and-wholesaling:pfof}{Amélioration de prix}~: 0.2 cent par action, 20
cents sur l'ordre. \omterm{def:m1:retail-flow-and-wholesaling:effective}{Demi-fourchette effective}~: $20.018 - 20.010 = 0.8$ cent.
Cinq minutes plus tard, le point médian vaut 20.013~: la \omterm{def:m1:retail-flow-and-wholesaling:effective}{demi-fourchette réalisée} est $20.018 -
20.013 = 0.5$ cent et l'impact sur le prix de 0.3 cent. Le revenu brut
du \omterm{def:m1:retail-flow-and-wholesaling:wholesaler}{wholesaler} sur cette exécution était de 50 cents, dont il verse une part au courtier
et couvre ses propres coûts.
\end{example}

\section{Pourquoi le flux non informé a de la valeur}

\begin{proposition}[Ce qu'un wholesaler peut se permettre de payer]\label{prop:m1:retail-flow-and-wholesaling:maxpay}
Soit $h$ la demi-fourchette cotée. Un flux dont les ordres sont informés avec
probabilité $p$, chaque ordre informé étant suivi d'un mouvement $J$, et qui
reçoit une \omterm{def:m1:retail-flow-and-wholesaling:pfof}{amélioration de prix} d'une fraction $\theta$ de $h$, procure une
\omterm{def:m1:retail-flow-and-wholesaling:effective}{demi-fourchette réalisée} espérée de $h(1-\theta) - pJ$. Un \omterm{def:m1:retail-flow-and-wholesaling:wholesaler}{wholesaler} ayant
des coûts d'exploitation $\kappa$ par action atteint le seuil de rentabilité en versant au courtier au plus
\[
\pi_{\max} \;=\; h(1-\theta) - pJ - \kappa .
\]
\end{proposition}
\begin{proof}
La \omterm{def:m1:retail-flow-and-wholesaling:effective}{demi-fourchette effective} $h(1-\theta)$, diminuée de l'impact espéré sur le prix $pJ$, est la
\omterm{def:m1:retail-flow-and-wholesaling:effective}{demi-fourchette réalisée}, d'après la \cref{prop:m1:retail-flow-and-wholesaling:identity}~;
retrancher les coûts.
\end{proof}

\begin{example}[Le même cent, deux flux]\label{ex:m1:retail-flow-and-wholesaling:twoflows}
Prenons $h = 1$ cent, $J = 3$ cents, $\kappa = 0.10$ cent. Sur une bourse,
où 35~\% des ordres frappant une cotation sont informés, un \omterm{def:m1:what-a-trading-firm-does:market-maker}{teneur de marché} gagne
$1 - 1.05 < 0$ avant coûts~: la fourchette cotée est, à l'équilibre, à peine
suffisante. Pour un flux de détail avec $p = 5\%$ et $\theta = 20\%$~: $\pi_{\max} =
0.80 - 0.15 - 0.10 = 0.55$ cent. Le \omterm{def:m1:retail-flow-and-wholesaling:wholesaler}{wholesaler} peut verser au courtier 0.25 cent
et garder 0.30 (voir la \cref{fig:m1:retail-flow-and-wholesaling:halfspread}). Rien
ici ne dépend de la vitesse ni du fait d'être plus malin que le client~: la valeur tient
entièrement à la \emph{segmentation} des ordres selon leur niveau d'information.
\end{example}

\begin{omfigure}
\begin{tikzpicture}
\begin{axis}[width=9cm,height=4.6cm,xbar,bar width=8pt,
  xlabel={cents par action, sur une demi-fourchette cotée de 1 cent},
  ytick=data,yticklabels={{Amélioration de prix pour le client},{Sélection adverse},{Coûts du wholesaler},{Paiement au courtier},{Marge du wholesaler}},
  yticklabel style={font=\small},tick label style={font=\small},label style={font=\small},
  y dir=reverse,xmin=0,xmax=0.4,xtick={0,0.1,0.2,0.3,0.4},enlarge y limits=0.15,grid=major,grid style={gray!25},
  xticklabel style={/pgf/number format/fixed,/pgf/number format/precision=1},
  nodes near coords,nodes near coords style={font=\footnotesize,/pgf/number format/fixed,/pgf/number format/precision=2},
  table/col sep=comma]
\addplot[fill=omDef!60,draw=omDef] table[x=cents,y=k]{figdata/markets-1/10-retail-flow-and-wholesaling/halfspread.csv};
\end{axis}
\end{tikzpicture}

\omcaption{Ce que devient un cent de demi-fourchette cotée lorsque l'ordre est
de détail. Les paramètres sont ceux de
l'\cref{ex:m1:retail-flow-and-wholesaling:twoflows} et sont donnés à titre illustratif. Données~:
le script du chapitre.}\label{fig:m1:retail-flow-and-wholesaling:halfspread}
\end{omfigure}

\begin{omfigure}
\begin{tikzpicture}
\begin{axis}[width=10.5cm,height=4.8cm,xmode=log,
  xlabel={temps après la transaction (secondes)},ylabel={demi-fourchette réalisée (cents)},
  tick label style={font=\small},label style={font=\small},
  xmin=0.1,xmax=300,ymin=-0.2,ymax=1.1,grid=major,grid style={gray!25},
  legend columns=2,legend style={font=\footnotesize,at={(0.5,-0.36)},anchor=north,draw=none,fill=none,/tikz/every even column/.append style={column sep=8pt}},
  table/col sep=comma]
\addplot[omDef,very thick,mark=*,mark size=1.4pt] table[x=horizon_s,y=retail_realised]{figdata/markets-1/10-retail-flow-and-wholesaling/markout.csv};
\addplot[omThm,very thick,mark=*,mark size=1.4pt] table[x=horizon_s,y=exchange_realised]{figdata/markets-1/10-retail-flow-and-wholesaling/markout.csv};
\legend{flux de détail (avec amélioration de 20~\%),flux frappant les cotations des bourses}
\end{axis}
\end{tikzpicture}

\omcaption{Ce que conserve le fournisseur de liquidité, selon le délai écoulé après la
transaction auquel il est mesuré (axe horizontal logarithmique). Les deux flux partent près de la
demi-fourchette cotée~; cinq minutes plus tard, le flux des bourses a repris la totalité de celle-ci
et le flux de détail en a conservé les deux tiers. Données~: la simulation du
tutoriel, 400\,000 transactions par courbe.}\label{fig:m1:retail-flow-and-wholesaling:markout}
\end{omfigure}

\begin{omfigure}
\begin{tikzpicture}
\begin{axis}[width=10.5cm,height=4.4cm,
  xlabel={part des ordres informés dans le flux (\%)},ylabel={$\pi_{\max}$ (cents par action)},
  tick label style={font=\small},label style={font=\small},
  xmin=0,xmax=30,ymin=-0.3,ymax=0.8,grid=major,grid style={gray!25},table/col sep=comma]
\addplot[omDef,very thick] table[x=informed_pct,y=max_payment]{figdata/markets-1/10-retail-flow-and-wholesaling/maxpay.csv};
\addplot[black,dashed] coordinates {(0,0) (30,0)};
\addplot[only marks,mark=*,omThm] coordinates {(5,0.55)};
\node[omThm,font=\footnotesize,anchor=south west,fill=white,inner sep=1.5pt] at (axis cs:5.6,0.57) {flux de détail de l'exemple};
\end{axis}
\end{tikzpicture}

\omcaption{Le montant maximal qu'un \omterm{def:m1:retail-flow-and-wholesaling:wholesaler}{wholesaler} peut payer pour un flux, selon le degré d'information
du flux, pour une demi-fourchette cotée d'un cent, une amélioration de 20~\%, un mouvement
informé de 3 cents et des coûts de 0.10 cent. Au-delà de 23~\% d'ordres informés, le flux
a un prix négatif. Données~: la \cref{prop:m1:retail-flow-and-wholesaling:maxpay},
calculées par le script du chapitre.}
\end{omfigure}

\begin{remark}[Qui paie]\label{rem:m1:retail-flow-and-wholesaling:whopays}
Le client de détail est mieux loti que sur une bourse, grâce à l'\omterm{def:m1:retail-flow-and-wholesaling:pfof}{amélioration de prix}
et à l'absence de commission. Le coût pèse ailleurs.
Les cotations des bourses sont fixées pour le flux qui atteint les bourses~; si les ordres les moins informés
en sont retirés, $p$ y augmente, et la fourchette aussi
(voir la \cref{rem:m1:us-equity-market-structure:free}). L'\omterm{def:m1:retail-flow-and-wholesaling:pfof}{amélioration de prix} est alors
mesurée par rapport à un indice de référence que le dispositif lui-même a élargi. Quelle
est l'ampleur de cet effet, et si une concurrence ordre par ordre servirait
mieux les clients que le wholesaling, a longtemps fait l'objet de débats devant la SEC,
sans réponse tranchée.
\end{remark}

\section{Quand le dispositif est mis à l'épreuve}

\begin{example}[Un cas de sanction]\label{ex:m1:retail-flow-and-wholesaling:case}
En décembre 2020, la SEC a poursuivi Robinhood Financial pour avoir induit en erreur
ses clients sur l'origine de ses revenus et pour avoir manqué à son devoir de meilleure
exécution~; la société a accepté de verser 65~millions de dollars, sans admettre ni nier
les constats. La décision relève qu'entre 2015 et la fin de 2018, les communications du courtier à ses clients
omettaient le \omterm{def:m1:retail-flow-and-wholesaling:pfof}{paiement pour flux d'ordres}, sa principale source
de revenus, et qu'il avait accepté des paiements anormalement élevés en échange d'une
\omterm{def:m1:retail-flow-and-wholesaling:pfof}{amélioration de prix} plus faible~: les exécutions de ses clients étaient de 34.1~millions de dollars moins bonnes
que chez les courtiers concurrents, \emph{même après} avoir tenu compte des
commissions qu'ils n'ont pas payées.
\end{example}

Le cas énonce le conflit avec précision~: la \cref{prop:m1:retail-flow-and-wholesaling:maxpay}
fixe le total, $h - pJ - \kappa$, et laisse son partage entre le
client ($\theta h$), le courtier ($\pi$) et le \omterm{def:m1:retail-flow-and-wholesaling:wholesaler}{wholesaler} à la négociation
entre ces deux derniers. Le client n'est pas à la table~; les
\omterm{def:m1:retail-flow-and-wholesaling:reports}{rapports Rule 605} et Rule 606, ainsi que l'obligation de meilleure exécution du courtier, sont là pour le représenter.

La semaine du 28 janvier 2021 a éprouvé la chaîne à un autre maillon~: non
pas la tarification mais la \emph{plomberie}. Les restrictions d'achat qui ont fait la une
provenaient du collatéral de la chambre de compensation (du \cref{ch:m1:clearing-and-settlement}),
et non des \omterm{def:m1:retail-flow-and-wholesaling:wholesaler}{wholesalers}~; le rapport du personnel de la SEC sur cet épisode
conclut que «~ce fut le sentiment positif, et non les rachats destinés à couvrir les positions vendeuses, qui a soutenu
la hausse du cours de l'action au centre de l'affaire pendant des semaines~».
Il s'agit du premier cas d'étude du \cref{ch:m1:stress-case-studies}.

\section{Tutoriel~: mesurer la qualité d'exécution}\label{tut:m1:retail-flow-and-wholesaling:execq}

\begin{tutorial}
\textbf{Objectif.} Calculer la fourchette effective, la fourchette réalisée et l'impact sur les prix
de deux flux d'ordres simulés et reproduire
la \cref{fig:m1:retail-flow-and-wholesaling:markout}. \textbf{État final~:} une
\omterm{def:m1:retail-flow-and-wholesaling:effective}{demi-fourchette réalisée} à cinq minutes de 0.65~cent pour le flux de détail et
environ zéro pour le flux des bourses.
\begin{tutsteps}
\item \textbf{Les flux.} Les ordres sont informés avec une probabilité de 5~\% (détail)
ou de 35~\% (bourses)~; un ordre informé est suivi d'un mouvement de 3~cents qui
se déploie sur une minute environ~; du bruit est ajouté à chaque horizon.
\omcode{code/markets-1/10-retail-flow-and-wholesaling/python/execq.py}{18}{38}{Deux flux qui ne diffèrent que par ce qu'ils savent et par l'amélioration de prix qui leur est accordée.}\label{lst:m1:retail-flow-and-wholesaling:flows}
\item \textbf{Les statistiques}~: trois moyennes, et l'identité qui les relie.
\omcode{code/markets-1/10-retail-flow-and-wholesaling/python/execq.py}{41}{46}{Fourchette effective, fourchette réalisée, impact sur les prix.}\label{lst:m1:retail-flow-and-wholesaling:stats}
\item \textbf{Calcul.} Les \omterm{def:m1:retail-flow-and-wholesaling:effective}{demi-fourchettes effectives} valent 0.80 et 1.00~cent. À 300
secondes, les impacts sur les prix valent 0.15 et 1.05~cents, soit les produits $pJ$~; les
\omterm{def:m1:retail-flow-and-wholesaling:effective}{demi-fourchettes réalisées} sont ce qui reste.
\item \textbf{Choisir l'horizon.} À 100 millisecondes, les deux flux paraissent
tout aussi rentables~; la différence n'apparaît qu'au bout de dizaines de secondes. Un
horizon plus court que le temps que met l'information pour atteindre le
prix surévalue chaque flux.
\end{tutsteps}
\textbf{Pour aller plus loin.} Donner aux ordres de détail portant sur un seul titre $p = 30\%$ (un
titre «~meme~» en pleine frénésie) et recalculer $\pi_{\max}$. Puis multiplier le bruit
par dix, et voir combien de transactions les moyennes exigent avant que les deux flux
puissent être distingués.
\end{tutorial}

\section{Construction~: le rapport de qualité d'exécution}\label{bld:m1:retail-flow-and-wholesaling:report}

\begin{build}
\bfield{Rôle} Chaque exécution de la société miniature, qu'elle fournisse
ou prenne de la liquidité, est notée de la même façon. Le rapport est le
propre tableau Rule 605 de la société et la matière première de l'analyse des coûts de
transaction du livre~10.
\bfield{Interface} \texttt{Fill(ts, symbol, side, quantity, price, bid, ask)}
avec la cotation à l'exécution~; un rappel \texttt{mid\_at(symbol, ts)}~;
\texttt{report(fills, horizons, by=(symbol, size\_bucket))} qui renvoie,
pour chaque groupe~: nombre d'exécutions et d'actions, \omterm{def:m1:retail-flow-and-wholesaling:effective}{demi-fourchettes effectives}, réalisées
et d'impact pondérées par les actions, exprimées en points de base, \omterm{def:m1:retail-flow-and-wholesaling:pfof}{amélioration de prix} en cents, et
la part des actions exécutées sur la cotation, à l'intérieur et à l'extérieur de celle-ci.
\bfield{Règles} Prix entiers~; l'identité de
la \cref{prop:m1:retail-flow-and-wholesaling:identity} doit être vérifiée exactement pour
chaque exécution. Une exécution avec une cotation croisée ou manquante est exclue et comptée
à part, jamais écartée en silence. Les tranches suivent les usages en matière de
taille d'ordre (1--99, 100--499, 500--1\,999, 2\,000--4\,999, 5\,000 et plus).
\bfield{Tests d'acceptation} \texttt{code/firm/execquality/tests/}~: l'exemple
à une seule exécution de ce chapitre~; l'identité sur des exécutions aléatoires~; la pondération
par les actions, et non par les exécutions~; les exclusions comptées.
\bfield{Pour aller plus loin} Ajouter la courbe de markout~: la fourchette réalisée pour une
liste d'horizons, avec une erreur-type par point.
\end{build}

\begin{omsources}
\item US Securities and Exchange Commission, \emph{SEC Charges Robinhood
Financial With Misleading Customers About Revenue Sources and Failing to
Satisfy Duty of Best Execution}, communiqué de presse 2020-321, 17 décembre 2020,
et ordonnance administrative 33-10906.
\item US Securities and Exchange Commission, \emph{Disclosure of Order
Execution Information}, texte d'adoption du 6 mars 2024~; \emph{Extension of
Compliance Date}, \emph{Federal Register}, 2 octobre 2025.
\item US Securities and Exchange Commission, \emph{Staff Report on Equity and
Options Market Structure Conditions in Early 2021}, octobre 2021.
\item Global Trading, ``Robinhood, Schwab led retail order flow payment
bonanza'', 2025 (compilation des rapports Rule 606).
\item R.~Huang and H.~Stoll, ``Dealer versus auction markets: a paired
comparison of execution costs on NASDAQ and the NYSE'', \emph{Journal of
Financial Economics} 41 (1996) --- fourchettes effectives et réalisées.
\item R.~Battalio, S.~Corwin and R.~Jennings, ``Can brokers have it all? On
the relation between make-take fees and limit order execution quality'',
\emph{Journal of Finance} 71 (2016).
\end{omsources}

\section{Exercices}

\begin{exercise}[$\star$]\label{exo:m1:retail-flow-and-wholesaling:1}
La cotation est $35.40 \times 35.44$. Un client vend 200 actions à 35.405.
Calculer l'\omterm{def:m1:retail-flow-and-wholesaling:pfof}{amélioration de prix} par action et en dollars, ainsi que la
\omterm{def:m1:retail-flow-and-wholesaling:effective}{demi-fourchette effective}.
\end{exercise}

\begin{exercise}[$\star$]\label{exo:m1:retail-flow-and-wholesaling:2}
Pour la vente de l'exercice précédent, le milieu une minute plus tard vaut 35.412.
Calculer la \omterm{def:m1:retail-flow-and-wholesaling:effective}{demi-fourchette réalisée} et l'impact sur le prix, puis vérifier
l'identité.
\end{exercise}

\begin{exercise}[$\star$]\label{exo:m1:retail-flow-and-wholesaling:3}
Un courtier reçoit 0.18 cent par action de \omterm{def:m1:retail-flow-and-wholesaling:pfof}{paiement pour flux d'ordres} sur 60 millions
d'actions par jour. Combien cela représente-t-il par an (252 jours)~? L'exprimer par client si
le courtier compte 12 millions de comptes approvisionnés.
\end{exercise}

\begin{exercise}[$\star\star$]\label{exo:m1:retail-flow-and-wholesaling:4}
Avec $h = 1.5$ cent, $J = 4$ cents, $\kappa = 0.12$ cent et une \omterm{def:m1:retail-flow-and-wholesaling:pfof}{amélioration de prix}
de 15~\% de $h$~: calculer $\pi_{\max}$ pour des flux avec $p = 4\%$ et
$p = 20\%$. À partir de quel $p$ le \omterm{def:m1:what-a-trading-firm-does:market-maker}{teneur de marché} cesse-t-il de faire une offre pour le flux~?
\end{exercise}

\begin{exercise}[$\star\star$]\label{exo:m1:retail-flow-and-wholesaling:5}
Le courtier A prélève 0.30 cent de paiement et ses clients reçoivent 10~\% de la
demi-fourchette en amélioration~; le courtier B prélève 0.10 cent et ses clients
reçoivent 30~\%. Avec $h = 1$ cent, de combien le client de B est-il mieux loti par action~? Si B
facture une commission de 1~\$ par ordre, pour quelles tailles d'ordre A est-il moins cher~?
\end{exercise}

\begin{exercise}[$\star\star$]\label{exo:m1:retail-flow-and-wholesaling:6}
Un \omterm{def:m1:retail-flow-and-wholesaling:wholesaler}{wholesaler} mesure des fourchettes réalisées de 0.70 cent à une seconde et de 0.20
cent à cinq minutes sur le flux d'un courtier, avec une \omterm{def:m1:retail-flow-and-wholesaling:effective}{demi-fourchette effective} de
0.75. Calculer l'impact sur le prix aux deux horizons. Que nous apprend la configuration
sur les clients de ce courtier, et que fera le \omterm{def:m1:retail-flow-and-wholesaling:wholesaler}{wholesaler}~?
\end{exercise}

\begin{exercise}[$\star\star\star$]\label{exo:m1:retail-flow-and-wholesaling:7}
\emph{Programmation.} Avec \texttt{simulate} et \texttt{stats}, estimer pour le
flux de détail l'erreur-type de la \omterm{def:m1:retail-flow-and-wholesaling:effective}{demi-fourchette réalisée} moyenne à 300
secondes à partir de $n = 10\,000$ transactions, puis en déduire combien de transactions sont nécessaires pour
distinguer $p = 5\%$ de $p = 8\%$ à deux erreurs-types.
\end{exercise}

\begin{exercise}[$\star\star\star$]\label{exo:m1:retail-flow-and-wholesaling:8}
\emph{Trouver la faille.} Le service marketing d'un \omterm{def:m1:retail-flow-and-wholesaling:wholesaler}{wholesaler} affirme~: «~Nous avons versé
1.2~milliard de dollars d'\omterm{def:m1:retail-flow-and-wholesaling:pfof}{amélioration de prix} aux investisseurs particuliers l'an dernier.~» Donner trois
raisons pour lesquelles ce chiffre, même calculé exactement, ne mesure pas le
bénéfice de ces investisseurs.
\end{exercise}

\section{Problème~: Tarification du flux de détail}

\begin{problem}[{Problème du week-end --- enchère pour le flux d'ordres d'un courtier}]\label{pb:m1:retail-flow-and-wholesaling:1}
Vous dirigez le desk de wholesaling d'une société de market making. Un courtier de détail qui achemine
40 millions d'actions par jour vous invite à soumettre une offre. Son flux porte sur des actions dont la
demi-fourchette cotée moyenne est de 1.2 cent. À partir d'un échantillon, vous estimez que 6~\%
de ses ordres sont informés, avec un mouvement ultérieur moyen de 3.5 cents. Vos coûts
s'élèvent à 0.11 cent par action.

\medskip\noindent\textbf{Partie I --- Économie du flux.}
\begin{enumerate}
\item Calculer l'impact attendu sur le prix par action.
\item Avec une \omterm{def:m1:retail-flow-and-wholesaling:pfof}{amélioration de prix} de 20~\% de la demi-fourchette, donner la
\omterm{def:m1:retail-flow-and-wholesaling:effective}{demi-fourchette effective} et la \omterm{def:m1:retail-flow-and-wholesaling:effective}{demi-fourchette réalisée}.
\item Calculer $\pi_{\max}$.
\item Donner votre revenu brut quotidien et votre marge quotidienne si vous payez 0.30 cent.
\item Et par an, sur 252 jours.
\end{enumerate}

\medskip\noindent\textbf{Partie II --- L'enchère.}
Le courtier classe les \omterm{def:m1:retail-flow-and-wholesaling:wholesaler}{wholesalers} selon l'\omterm{def:m1:retail-flow-and-wholesaling:pfof}{amélioration de prix} délivrée et fixe
lui-même le paiement à 0.30 cent pour tous.
\begin{enumerate}[resume]
\item Quelle est la plus grande fraction d'amélioration $\theta$ à laquelle vous atteignez
le seuil de rentabilité~?
\item Un concurrent a des coûts de 0.06 cent. Quel $\theta$ peut-il offrir~?
\item Vous pouvez couvrir l'inventaire résiduel à moindre coût que lui, ce qui réduit
votre impact effectif sur le prix de 0.08 cent. Qui l'emporte désormais~?
\item Pourquoi le courtier achemine-t-il ses ordres vers plusieurs \omterm{def:m1:retail-flow-and-wholesaling:wholesaler}{wholesalers} à la fois et redistribue-t-il
les actions chaque mois~?
\end{enumerate}

\medskip\noindent\textbf{Partie III --- Le flux évolue.}
\begin{enumerate}[resume]
\item Pendant une frénésie spéculative, la part d'ordres informés dans vingt
actions passe à 25~\% et le mouvement à 6 cents. Calculer la
\omterm{def:m1:retail-flow-and-wholesaling:effective}{demi-fourchette réalisée} dans ces actions pour $\theta = 20\%$ et des fourchettes inchangées.
\item Ces actions représentent 15~\% du volume du courtier. Donner votre nouvelle marge
quotidienne, en payant 0.30 cent.
\item Les fourchettes de ces actions s'élargissent à une demi-fourchette de 3 cents. Refaire la
question 10.
\item Énumérer trois actions possibles, dans le cadre de votre accord avec le courtier,
lorsqu'un segment du flux devient toxique.
\end{enumerate}

\medskip\noindent\textbf{Partie IV --- Les autres parties.}
\begin{enumerate}[resume]
\item Que gagne le courtier par an grâce à votre paiement~?
\item Que reçoivent ses clients par an en \omterm{def:m1:retail-flow-and-wholesaling:pfof}{amélioration de prix}~?
\item Le courtier pourrait plutôt envoyer tous ses ordres à une bourse qui facture
des frais preneurs de 0.30 cent. Que paieraient les clients par rapport au milieu, et
que paierait le courtier~?
\item Comparer le coût total des clients dans les deux dispositifs.
Qu'a-t-on supposé sur la fourchette cotée de la bourse~?
\item Dans le cadre de la \cref{rem:m1:retail-flow-and-wholesaling:whopays}, comment cette
hypothèse pourrait-elle échouer si tout le flux de détail migrait vers les bourses~?
\item Énoncer le \emph{résultat nommé}~: le paiement maximal par action auquel
vous atteignez le seuil de rentabilité sur ce flux, en cents.
\item En une phrase~: qu'achète exactement le \omterm{def:m1:retail-flow-and-wholesaling:wholesaler}{wholesaler}~?
\end{enumerate}
\end{problem}

\section{Questions d'entretien}

\begin{interviewq}[$\star$ \iqroles{trader, chercheur, développeur}]\label{iq:m1:retail-flow-and-wholesaling:1}
Qu'est-ce que le \omterm{def:m1:retail-flow-and-wholesaling:pfof}{paiement pour flux d'ordres}, et pourquoi un \omterm{def:m1:what-a-trading-firm-does:market-maker}{teneur de marché} paierait-il pour recevoir des ordres~?
\end{interviewq}

\begin{interviewq}[$\star$ \iqroles{trader, chercheur}]\label{iq:m1:retail-flow-and-wholesaling:2}
Définissez la \omterm{def:m1:retail-flow-and-wholesaling:effective}{demi-fourchette effective} et la \omterm{def:m1:retail-flow-and-wholesaling:effective}{demi-fourchette réalisée}. Laquelle mesure le
coût du client et laquelle mesure le revenu du \omterm{def:m1:what-a-trading-firm-does:market-maker}{teneur de marché}~?
\end{interviewq}

\begin{interviewq}[$\star\star$ \iqroles{chercheur, ingénieur ML}]\label{iq:m1:retail-flow-and-wholesaling:3}
On vous fournit une année d'exécutions d'un courtier, avec les cotations correspondantes. Comment décidez-vous
si son flux est «~toxique~», et à quel horizon mesurez-vous~?
\end{interviewq}

\begin{interviewq}[$\star\star$ \iqroles{trader, chercheur}]\label{iq:m1:retail-flow-and-wholesaling:4}
Un ordre de détail et un ordre enfant d'un algorithme institutionnel portent
tous deux sur 100 actions au marché. Pourquoi un \omterm{def:m1:what-a-trading-firm-does:market-maker}{teneur de marché} les traite-t-il différemment, et comment
peut-il les distinguer sur une bourse anonyme~?
\end{interviewq}

\begin{interviewq}[$\star\star$ \iqroles{chercheur, banque}]\label{iq:m1:retail-flow-and-wholesaling:5}
Défendez la thèse selon laquelle le \omterm{def:m1:retail-flow-and-wholesaling:pfof}{paiement pour flux d'ordres} est bénéfique aux investisseurs de détail, puis
celle selon laquelle il nuit aux marchés.
\end{interviewq}

\begin{interviewq}[$\star\star\star$ \iqroles{chercheur, trader}]\label{iq:m1:retail-flow-and-wholesaling:6}
Sur le flux d'un courtier, la \omterm{def:m1:retail-flow-and-wholesaling:effective}{demi-fourchette réalisée} de votre desk de wholesaling a été divisée par deux
en six mois, alors que sa \omterm{def:m1:retail-flow-and-wholesaling:effective}{demi-fourchette effective} est restée inchangée. Donnez trois hypothèses
et les données qui permettraient de les départager.
\end{interviewq}
